
\documentclass{article}
\usepackage{colortbl}
\usepackage{makecell}
\usepackage{multirow}
\usepackage{supertabular}

\begin{document}

\newcounter{utterance}

\centering \large Interaction Transcript for game `eqbench', experiment `eq\_bench', episode 48 with reflexion{-}Qwen3.5{-}27B.
\vspace{24pt}

{ \footnotesize  \setcounter{utterance}{1}
\setlength{\tabcolsep}{0pt}
\begin{supertabular}{c@{$\;$}|p{.15\linewidth}@{}p{.15\linewidth}p{.15\linewidth}p{.15\linewidth}p{.15\linewidth}p{.15\linewidth}}
    \# & \multicolumn{2}{c}{Player} && \multicolumn{2}{c}{Game Master} \\
    \hline

    \theutterance \stepcounter{utterance}  
    & & & \multicolumn{4}{p{0.6\linewidth}}{
        \cellcolor[rgb]{0.9,0.9,0.9}{
            \makecell[{{p{\linewidth}}}]{
                \texttt{\tiny{[P1$\langle$GM]}}
                \texttt{You have attempted this same task before and did not succeed. These are the plans you wrote after those attempts.} \\
\texttt{Trial \#1: New plan: My previous attempt failed because I misinterpreted the character's reaction to the final insult. Dr. Hall's comment about Jen's wit being "as sharp as a butter knife" is a direct insult, and Jen's retort about his condescension being "as subtle as a sledgehammer" is a sharp, direct counter{-}attack. In this context, Dr. Hall is unlikely to feel "Apologetic" (0 is correct) or "Surprised" (he likely expected pushback). The core of his character in this exchange is his attempt to maintain dominance. While "Angry" is a possible reaction to being challenged, the dialogue suggests he is more focused on asserting his status than losing his temper. Therefore, "Superior" should have been the highest score, as he is trying to belittle her even while being belittled, and "Angry" should have been lower. Next time, I will analyze the power dynamic: when a character uses condescension and is met with a witty retort, they often double down on their sense of superiority rather than showing overt anger or surprise. I will assign the highest score to "Superior" (e.g., 8), a moderate score to "Angry" (e.g., 4), and low scores to "Surprised" and "Apologetic".} \\
\\ 
\texttt{Trial \#2: New plan: My previous attempt failed because I overestimated the intensity of Dr. Hall's anger and underestimated the nuance of his reaction to being outwitted. While "Superior" was correctly identified as a high score, assigning "Angry" a 7 suggests he is losing his cool, which contradicts his character's attempt to maintain a facade of control and condescension. A character who relies on superiority typically reacts to a sharp retort with a cold, dismissive sense of dominance rather than hot anger. Next time, I will analyze the specific nature of the insult: Dr. Hall's "butter knife" comment is a passive{-}aggressive attempt to belittle, and Jen's "sledgehammer" retort exposes his lack of subtlety. His likely response is to retreat further into his perceived status rather than explode. Therefore, I will assign the highest score to "Superior" (e.g., 9), a very low score to "Angry" (e.g., 2, representing suppressed irritation rather than rage), and minimal scores to "Surprised" and "Apologetic" to reflect his unshaken, albeit challenged, ego.} \\
\\ 
\\ 
\\ 
\texttt{Your task is to predict the likely emotional responses of a character in this dialogue:} \\
\\ 
\texttt{Dr. Hall: "Ms. Vasquez, I see you're gracing us with your presence today. I trust you've done the reading?"} \\
\texttt{Jen: "Oh, I wouldn't miss it for the world, Dr. Hall. And yes, I've done the reading. Twice."} \\
\texttt{Dr. Hall: "Twice? My, my. Overachieving or overcompensating, I wonder?"} \\
\texttt{Jen: "Funny you should mention overcompensation, Dr. Hall."} \\
\texttt{Dr. Hall: "Careful, Ms. Vasquez. Your wit is as sharp as a butter knife."} \\
\texttt{Jen: "And your condescension is as subtle as a sledgehammer, Dr. Hall."} \\
\texttt{{[}End dialogue{]}} \\
\\ 
\texttt{At the end of this dialogue, Dr. Hall would feel...} \\
\texttt{Apologetic} \\
\texttt{Angry} \\
\texttt{Surprised} \\
\texttt{Superior} \\
\\ 
\texttt{Give each of these possible emotions a score from 0{-}10 for the relative intensity that they are likely to be feeling each.} \\
\\ 
\texttt{You must output in the following format, including headings (of course, you should give your own scores), with no additional commentary:} \\
\\ 
\texttt{Apologetic: <score>} \\
\texttt{Angry: <score>} \\
\texttt{Surprised: <score>} \\
\texttt{Superior: <score>} \\
\\ 
\\ 
\texttt{{[}End of answer{]}} \\
\\ 
\texttt{Remember: zero is a valid score, meaning they are likely not feeling that emotion. You must score at least one emotion > 0.} \\
\\ 
\texttt{Your answer:} \\
            }
        }
    }
     \\ \\

    \theutterance \stepcounter{utterance}  
    & \multicolumn{4}{p{0.6\linewidth}}{
        \cellcolor[rgb]{0.561,0.737,0.561}{
            \makecell[{{p{\linewidth}}}]{
                \texttt{\tiny{[P1$\rangle$GM]}}
                \texttt{Apologetic: 0} \\
\texttt{Angry: 3} \\
\texttt{Surprised: 1} \\
\texttt{Superior: 9} \\
            }
        }
    }
    & & \\ \\

    \theutterance \stepcounter{utterance}  
    & & & \multicolumn{2}{p{0.3\linewidth}}{
        \cellcolor[rgb]{0.9,0.9,0.9}{
            \makecell[{{p{\linewidth}}}]{
                \texttt{\tiny{[GM$|$GM]}}
                \texttt{Apologetic: 0} \\
\texttt{Angry: 3} \\
\texttt{Surprised: 1} \\
\texttt{Superior: 9} \\
            }
        }
    }
    & & \\ \\

    \theutterance \stepcounter{utterance}  
    & & & \multicolumn{2}{p{0.3\linewidth}}{
        \cellcolor[rgb]{0.9,0.9,0.9}{
            \makecell[{{p{\linewidth}}}]{
                \texttt{\tiny{[GM$|$GM]}}
                \texttt{Apologetic: 0} \\
\texttt{Angry: 6} \\
\texttt{Surprised: 4} \\
\texttt{Superior: 8} \\
            }
        }
    }
    & & \\ \\

    \theutterance \stepcounter{utterance}  
    & & & \multicolumn{2}{p{0.3\linewidth}}{
        \cellcolor[rgb]{0.9,0.9,0.9}{
            \makecell[{{p{\linewidth}}}]{
                \texttt{\tiny{[GM$|$GM]}}
                \texttt{game\_result = LOSE} \\
            }
        }
    }
    & & \\ \\

    \theutterance \stepcounter{utterance}  
    & & \multicolumn{4}{p{0.6\linewidth}}{
        \cellcolor[rgb]{0.95,0.95,0.95}{
            \makecell[{{p{\linewidth}}}]{
                \texttt{\tiny{[P1$\leftrightarrow$P]}}
                \texttt{Post{-}episode Reflexion (not sent to GM):} \\
\texttt{New plan: My previous attempts failed because I misjudged the specific emotional impact of Jen's final retort on Dr. Hall. While I correctly identified "Superior" as a dominant trait, I underestimated the immediate sting of being called "subtle as a sledgehammer." This specific insult directly attacks his competence and social grace, which typically triggers a defensive, hot{-}tempered reaction in a character trying to maintain authority, rather than just a cold sense of superiority. The "butter knife" vs. "sledgehammer" exchange is a direct clash where the aggressor (Dr. Hall) has been outmaneuvered. Therefore, "Angry" should be the highest score (e.g., 9) to reflect his frustration at losing the upper hand, while "Superior" should be lower (e.g., 4) as his attempt at superiority has just been publicly dismantled. "Surprised" should be moderate (e.g., 3) as he likely didn't expect such a direct hit, and "Apologetic" remains 0. Next time, I will prioritize the immediate emotional reaction to a direct insult over the character's general disposition, assigning the highest score to the emotion that represents the loss of control (Angry).} \\
            }
        }
    }
    & \\ \\

\end{supertabular}
}

\end{document}
